% from ega2-1-fr.tex
\oldpage[II]{43}
On notera qu'il y a en général des morphismes de $\operatorname{Proj}(S)$
dans $\operatorname{Proj}(S')$ qui ne sont pas affines, et ne proviennent
donc pas d'homomorphismes d'anneaux gradués $S'\to S$~; un exemple est le
morphisme structural $\operatorname{Proj}(S)\to\operatorname{Spec}(A)$ quand
$A$ est un corps ($\operatorname{Spec}(A)$ étant identifié à
$\operatorname{Proj}(A[T])$ (cf. (\hyperref[II.3.1.7-fr]{3.1.7})))~; cela
résulte en effet de (\hyperref[I.2.3.2-fr]{I, 2.3.2}).

\begin{env}[2.8.4]
\label{II.2.8.4-fr}
Soient $S''$ un troisième anneau gradué à degrés positifs,
$\varphi':S''\to S'$ un homomorphisme d'anneaux gradués, et posons
$\varphi''=\varphi\circ\varphi'$. En raison de
(\hyperref[II.2.8.1.1-fr]{2.8.1.1}) et de la formule
${}^a\varphi''=({}^a\varphi')\circ({}^a\varphi)$ on vérifie aussitôt que l'on
a $G(\varphi'')\subset G(\varphi)$ et que, si $\Phi$, $\Phi'$ et $\Phi''$
sont les morphismes associés à $\varphi$, $\varphi'$ et $\varphi''$, on a
$\Phi''=\Phi'\circ(\Phi|G(\varphi''))$.
\end{env}

\begin{env}[2.8.5]
\label{II.2.8.5-fr}
Supposons que $S$ (resp. $S'$) soit une $A$-algèbre graduée (resp. une
$A'$-algèbre graduée), et soit $\psi:A'\to A$ un homomorphisme d'anneaux tel
que le diagramme
\[
\label{II.2.8.5.1-fr}
  \xymatrix{
    A'\ar[r]^{\psi}\ar[d] & A\ar[d]\\
    S'\ar[r]_{\varphi} & S
  }
\tag{2.8.5.1}
\]
soit commutatif. On peut alors considérer $G(\varphi)$ et
$\operatorname{Proj}(S')$ comme des schémas au-dessus de
$\operatorname{Spec}(A)$ et $\operatorname{Spec}(A')$ respectivement~; si
$\Phi$ et $\Psi$ sont les morphismes associés respectivement à $\varphi$ et
$\psi$, le diagramme
\[
  \xymatrix{
    G(\varphi)\ar[r]^{\Phi}\ar[d] & \operatorname{Proj}(S')\ar[d]\\
    \operatorname{Spec}(A)\ar[r]_{\Psi} & \operatorname{Spec}(A')
  }
\]
est alors commutatif~: il suffit en effet de le démontrer pour la restriction
de $\Phi$ à $D_+(f)$, où $f=\varphi(f')$, $f'$ homogène dans $S'_+$~; et
alors cela résulte de la commutativité du diagramme
\[
  \xymatrix{
    A'\ar[r]^{\psi}\ar[d] & A\ar[d]\\
    S'_{(f')}\ar[r]_{\varphi_{(f)}} & S_{(f)}
  }
\]
\end{env}

\begin{env}[2.8.6]
\label{II.2.8.6-fr}
Soit maintenant $M$ un $S$-module gradué, et considérons le $S'$-module
$M_{[\varphi]}$, qui est évidemment gradué. Soit $f'$ homogène dans $S'_+$,
et soit $f=\varphi(f')$~; on sait
(\hyperref[0.1.5.2-fr]{0, 1.5.2}) qu'il y a un isomorphisme canonique
$(M_{[\varphi]})_{f'}\xrightarrow{\sim}(M_f)_{[\varphi_f]}$, et il est
immédiat que cet isomorphisme conserve les degrés, donc il y a un isomorphisme
canonique
$(M_{[\varphi]})_{(f')}\xrightarrow{\sim}
(M_{(f)})_{[\varphi_{(f)}]}$. Il correspond canoniquement à cet isomorphisme
un isomorphisme de faisceaux
$\widetilde{M_{[\varphi]}}|D_+(f')
\xrightarrow{\sim}(\Phi_f)_*(\widetilde{M}|D_+(f))$
(\hyperref[II.2.5.2-fr]{2.5.2} et
\hyperref[I.1.6.3-fr]{I, 1.6.3}). En outre,

